% ============================================================================
% Chapter 1: Anchor Junction as Topological Ground State
% ============================================================================
% Source: 04b_proton_anchor.tex, 04c_routeB_z6_steiner.tex
% Status: FILLED from SoT
% Contamination check: PASSED (no SM terms)
% ============================================================================

\chapter{Anchor Junction as Topological Ground State}
\label{ch:anchor}

\source{04b\_proton\_anchor.tex}

\begin{abstract}
Why is the anchor junction stable? Experimental bounds exceed $\tau_p > 10^{34}$ years, yet no
conventional symmetry principle mandates this. In the EDC framework, anchor stability
emerges from topology: the anchor junction corresponds to a \emph{topological anchor}---a locally
minimizing Y-junction configuration of 5D worldsheets that cannot decay via continuous
deformations.
\end{abstract}

% ----------------------------------------------------------------------------
% CHAPTER SPINE
% ----------------------------------------------------------------------------
\paragraph{Chapter Spine.}
\textbf{Purpose:} Establish the anchor junction ($\Zsix$ ground state) as the
topological minimum of 5D brane configurations, explaining its exceptional
stability.
\textbf{Inputs:} (i)~Nambu--Goto action for worldsheets~[P], (ii)~classical
Steiner theorem~[M], (iii)~$\pi_1$ homotopy obstruction~[M].
\textbf{Outputs:} (i)~Anchor junction is local energy minimum within Y-junction
sector~[Dc], (ii)~$120°$ Steiner angles are geometrically mandated~[Dc],
(iii)~topological protection prevents decay to trivial sector~[Dc].
\textbf{Dependencies:} None (foundational chapter).
\textbf{Forward links:} Supplies $\Zsix$ symmetry to Ch.~\ref{ch:junction} and
ground-state reference for metastable decay in Ch.~\ref{ch:metastable}.

% ============================================================================
\section{The Stability Puzzle}
\label{sec:anchor:puzzle}

The anchor junction is the lightest \Junction{} configuration. Experimental searches constrain
its lifetime:
\begin{equation}
    \tau_p > 10^{34} \, \text{years} \tagBL
\end{equation}

\noindent
\textbf{The puzzle:} Why should this configuration be stable at all? What prevents it
from decaying to lighter states?

In EDC, the answer is topological. The anchor junction is not merely ``long-lived by accident''
but is a \emph{local minimum} of a 5D energy functional, protected by topology from
decay to the trivial vacuum.

% ============================================================================
\section{The Y-Junction Model}
\label{sec:anchor:yjunction}

\subsection{Postulate: Anchor as Topological Anchor}

\begin{derivationbox}[title=Postulate: Anchor Stability Principle]
    The anchor junction Y-junction configuration represents a \emph{local minimum} of an
    appropriate 5D energy functional $\mathcal{E}[\Psi]$ under thick-brane boundary
    conditions. The anchor junction is a \textbf{topological anchor} that stabilizes the
    brane--observer interface. \tagP
\end{derivationbox}

\subsection{5D Energy Functional}

We assume the existence of an effective 5D energy functional \tagP:
\begin{equation}
    \label{eq:5d_energy_functional}
    \mathcal{E}[\Psi] \;=\; \mathcal{E}_{\mathrm{bulk}}[\Psi]
        \;+\; \mathcal{E}_{\mathrm{brane}}[\Psi]
        \;+\; \mathcal{E}_{\mathrm{BC}}[\Psi]
\end{equation}
whose configuration space $\mathcal{C}$ partitions into topologically distinct sectors.
The anchor junction occupies a sector $\mathcal{C}_Y$ (Y-junction configurations) separated
from the trivial sector by an energy barrier.

\subsection{Functional Role}

In the Book II derivation pipeline, the anchor junction serves as the \textbf{benchmark}
for energy accounting. When the metastable junction decays, the anchor is the ground-state
endpoint---the stable configuration to which the system relaxes. Without a stable
anchor junction, the entire ledger-closure mechanism would lack a fixed reference point.

% ============================================================================
\section{Route A: Topological Proof of Stability}
\label{sec:anchor:routeA}

\source{04b\_proton\_anchor.tex, \S Route A}

The following chain establishes anchor stability using only topological protection,
the Nambu--Goto action, and the classical Steiner theorem. This derivation is
\emph{independent} of crystallographic structure.

\subsection{Lemma 1: Topological Protection}

\begin{derivationbox}[title=Lemma 1: Topological Protection of Worldsheets]
    Let $\Sigma \subset M_5$ be a 2D worldsheet with boundary $\partial\Sigma$
    fixed on the brane $\mathcal{B}$. If $\pi_1(\mathcal{B} \setminus \partial\Sigma) \neq 0$,
    then $\Sigma$ cannot be continuously contracted to a point while keeping its
    boundary fixed.

    \textbf{Consequence:} Worldsheets with non-trivial winding are topologically
    stable---they cannot decay via smooth deformations. \tagDer
\end{derivationbox}

\begin{proof}
By hypothesis, the boundary curve $\partial\Sigma$ represents a non-trivial element
of $\pi_1(\mathcal{B} \setminus \{\text{defect locus}\})$. Any continuous deformation
of $\Sigma$ that contracts it to a point would require $\partial\Sigma$ to become
contractible, contradicting the non-triviality. This is a standard result in algebraic
topology.
\end{proof}

\subsection{Lemma 2: Nambu--Goto Energy}

\begin{derivationbox}[title=Lemma 2: Nambu--Goto Energy for Frozen Configurations]
    For a worldsheet $\Sigma$ with constant tension $\tau$, the Nambu--Goto action
    in the frozen (static) limit gives:
    \begin{equation}
        \label{eq:nambu_goto_energy}
        E[\Sigma] \;=\; \tau \cdot \mathrm{Length}(\Sigma)
    \end{equation}
    where $\mathrm{Length}(\Sigma)$ is the total worldsheet length projected onto
    the spatial brane. \tagDer
\end{derivationbox}

\begin{proof}
The Nambu--Goto action is $S_{\mathrm{NG}} = \tau \int d^2\sigma \sqrt{-\det(h_{ab})}$,
where $h_{ab}$ is the induced metric. In the static gauge where the worldsheet is
time-independent, the temporal integral factorizes:
$S_{\mathrm{NG}} = \tau \cdot T \cdot \int d\sigma \sqrt{g_{\sigma\sigma}}$.
The spatial integral is the arc length. The energy $E = S_{\mathrm{NG}}/T = \tau L$.
\end{proof}

\subsection{Theorem: Steiner Optimality}

\begin{resultbox}[title=Theorem: Steiner Optimality for Y-Junctions]
    Let $\{A, B, C\}$ be three fixed points on a 2D surface. Among all connected
    networks joining these three points, the \textbf{Steiner tree}---meeting at a
    single interior vertex with $120°$ angles---minimizes total length, provided
    all three edge tensions are equal. \tagM
\end{resultbox}

\begin{proof}[Reference]
This is the classical Steiner problem, solved by Fermat, Torricelli, and Steiner.
The $120°$ condition follows from force balance: three equal-tension lines meeting
at a point are in equilibrium if and only if the angles between them are $120°$.
See Gilbert \& Pollak (1968).
\end{proof}

\subsection{Corollary: Anchor as Minimum}

\begin{resultbox}[title=Corollary: Anchor as Topological-Geometric Minimum]
    The anchor junction Y-junction---three worldsheets meeting at $120°$ angles---is a
    \emph{local minimum} of the 5D energy functional within its topological sector:

    \begin{enumerate}
        \item \textbf{Topological protection} (Lemma 1): The configuration cannot
              decay to the trivial sector.
        \item \textbf{Length minimization} (Lemma 2 + Theorem): Among all Y-junction
              configurations with the same boundary conditions, the $120°$ Steiner
              configuration minimizes $E = \tau L$.
        \item \textbf{Stability}: The second variation $\delta^2 E > 0$ for
              perturbations preserving the topology (positive Hessian).
    \end{enumerate}
    \tagDc
\end{resultbox}

\begin{proof}
Combine Lemmas 1 and 2 with the Steiner theorem. The topological sector $\mathcal{C}_Y$
is preserved by continuous deformations (Lemma 1). Within $\mathcal{C}_Y$, the energy
is $E = \tau L$ (Lemma 2), so minimizing energy is equivalent to minimizing length.
The Steiner theorem identifies the $120°$ Y-junction as the unique length minimizer.
Positive Hessian follows from the strict local minimum property of Steiner configurations.
\end{proof}

% ============================================================================
\section{Route B: Crystallographic Derivation}
\label{sec:anchor:routeB}

\source{04c\_routeB\_z6\_steiner.tex}

An independent derivation reaches the same $120°$ result via the $\Zsix$
crystallization program.

\subsection{Postulate P2: Worldsheet Interactions}

\begin{derivationbox}[title=Postulate P2: Worldsheet Interactions]
    Worldsheets (topological defects) in the thick-brane have:
    \begin{enumerate}
        \item \textbf{Short-range repulsion:} Core overlap creates energy divergence
              ($V_{\mathrm{core}} \to +\infty$ as $d \to 0$)
        \item \textbf{Long-range confinement:} Logarithmic or linear growth
              ($V_{\mathrm{lin}} \to +\infty$ as $d \to \infty$)
    \end{enumerate}
    The combined potential $V(d) = V_{\mathrm{core}}(d) + V_{\mathrm{lin}}(d)$ has
    a minimum at some characteristic distance $d_0 > 0$. \tagP
\end{derivationbox}

\subsection{Kepler--Hales Lemma}

\begin{derivationbox}[title=Lemma B1: Kepler--Hales Optimal Packing]
    The densest packing of equal circles in 2D is the hexagonal lattice, with
    packing fraction $\pi/(2\sqrt{3}) \approx 90.69\%$. \tagM
\end{derivationbox}

\subsection{Crystallization Chain}

The Route B derivation proceeds:

\begin{enumerate}
    \item \textbf{[P]} Postulate P2: Worldsheets have repulsion + confinement with
          minimum at $d_0$
    \item \textbf{[M]} Kepler--Hales: Optimal 2D packing is hexagonal
    \item \textbf{[Dc]} Crystallization: Worldsheets form hexagonal lattice
    \item \textbf{[Dc]} $\Zsix$ symmetry: The $\Zthree \subset \Zsix$ subgroup at
          Y-junctions implies equal tensions along all three branches
    \item \textbf{[M]} Steiner theorem: Equal tensions $\Rightarrow$ $120°$ angles
\end{enumerate}

\begin{resultbox}[title=Corollary: Anchor $120°$ via Route B]
    The anchor junction Y-junction exhibits $120°$ Steiner geometry as a consequence of
    $\Zsix$ crystallization on the brane. \tagDc
\end{resultbox}

% ============================================================================
\section{Convergence of Routes}
\label{sec:anchor:convergence}

\begin{resultbox}[title=Convergence Statement]
    \textbf{Route A:} Topology + Nambu--Goto + Steiner
    \[
        \pi_1 \text{ protection [M]+[P]} \to E = \tau L \text{ [Der]}
        \to \text{Steiner } 120° \text{ [M]} \to \text{Anchor minimum [Dc]}
    \]

    \textbf{Route B:} P2 + Crystallization + Steiner
    \[
        \text{P2 [P]} \to \text{Kepler--Hales [M]} \to \text{Hex lattice [Dc]}
        \to \text{Equal tensions [Dc]} \to \text{Steiner } 120° \text{ [M]}
    \]

    \textbf{Common terminus:} Both routes use the Steiner theorem as the final
    mathematical step. The $120°$ result is \emph{overdetermined}---two independent
    derivation chains converge on the same geometry.
\end{resultbox}

% ============================================================================
\section{Role of Boundary Conditions}
\label{sec:anchor:BC}

\begin{edcopenbox}[title=Clarification: Boundary Conditions]
    \textbf{What BC do:}
    \begin{itemize}
        \item Provide the scale $\bdelta$ (brane thickness)
        \item Set the mode spectrum of fluctuations
    \end{itemize}

    \textbf{What BC do NOT do:}
    \begin{itemize}
        \item BC do \emph{not} create the attractive interaction that produces
              the energy minimum
        \item Linearized BC analysis gives $V'_{\mathrm{lin}}(d) > 0$ for all BC
              (Neumann, Robin, Dirichlet)
    \end{itemize}

    The minimum arises from the balance:
    \begin{itemize}
        \item Short range: Topological core repulsion ($V_{\mathrm{core}} \to +\infty$
              as $d \to 0$)
        \item Long range: Logarithmic confinement ($V_{\mathrm{lin}} \to +\infty$
              as $d \to \infty$)
        \item Continuity: A minimum exists at some $d_0 > 0$
    \end{itemize}
\end{edcopenbox}

% ============================================================================
\section{Epistemic Status}
\label{sec:anchor:epistemic}

\begin{center}
\begin{tabular}{llll}
\toprule
\textbf{Claim} & \textbf{Status} & \textbf{Reference} & \textbf{Source} \\
\midrule
Topological protection & [M]+[P] & Lemma 1 & $\pi_1$ obstruction \\
$E = \tau L$ (Nambu--Goto) & [Der] & Lemma 2 & Action variation \\
$120°$ Steiner angles & [M] & Theorem & Classical geometry \\
Anchor junction is Y-junction minimum & [Dc] & Corollary & Proof chain \\
Positive Hessian & [Dc] & Corollary & Steiner stability \\
\midrule
$\Zsix$ crystallization & [Dc] & Route B & P2 + Kepler--Hales \\
BC role clarified & [Der] & \S\ref{sec:anchor:BC} & BC audit \\
\midrule
Explicit $\mathcal{E}[\Psi]$ form & \tagOpen & --- & Future work \\
Barrier height (metastability) & \tagOpen & --- & Future work \\
\bottomrule
\end{tabular}
\end{center}

% ============================================================================
\section{Falsifiability}
\label{sec:anchor:falsify}

\begin{edcopenbox}[title=Falsifiability Hooks]
    \begin{itemize}
        \item If anchor decay is observed at rates inconsistent with a topologically
              protected minimum, the anchor mechanism fails.
        \item If the Y-junction configuration cannot be realized as a stationary
              point of any reasonable 5D energy functional, the structural claim
              is falsified.
        \item If the second variation $\delta^2\mathcal{E}$ has negative eigenvalues
              (unstable directions), the local-minimum claim fails.
        \item If \Junction{} states can be destabilized by small perturbations without
              violating conservation laws, the topological protection is illusory.
    \end{itemize}
\end{edcopenbox}

% ============================================================================
% Observer-side projection
% ============================================================================

\begin{observerbox}
Observer-side projection note: quoted terms are measurement labels for the
3D/4D projection of the 5D objects defined in this chapter; they do not
imply any conventional mechanism.

\begin{itemize}
    \item \AnchorJunction{} $\leftrightarrow$ ``proton'' (projection label)
\end{itemize}

What the observer records: a stable, positively-charged junction state that
serves as the ground-state reference for energy accounting. The measured
mass and charge are projection-level observables.
\end{observerbox}

% ============================================================================
\section{Summary}
\label{sec:anchor:summary}

\begin{resultbox}
    The anchor junction is the unique topological ground state of the 5D brane configuration,
    characterized by:
    \begin{itemize}
        \item $\Zsix$ junction symmetry (with $\Zthree$ at Y-vertices)
        \item $120°$ Steiner angles (from length/energy minimization)
        \item Topological protection (cannot decay to trivial sector)
        \item Positive Hessian (stable local minimum)
    \end{itemize}

    Anchor stability follows from topology, not from imposed symmetries.
\end{resultbox}

% ----------------------------------------------------------------------------
% BRIDGE TO NEXT CHAPTER
% ----------------------------------------------------------------------------
\paragraph{Bridge to Chapter~\ref{ch:junction}.}
The $\Zsix$ symmetry established here is not merely a label---it is a
\emph{generative} symmetry whose subgroup structure ($\Zthree$, $\Ztwo$)
governs the spectrum of junction states. Chapter~\ref{ch:junction} unpacks
this subgroup chain and introduces the $\Msix$ coordination lattice that
emerges when multiple junctions tile the brane. Understanding the lattice
is prerequisite to the metastable junction (Ch.~\ref{ch:metastable}) and
the pinning mechanism (Ch.~\ref{ch:sigma_to_K}).
